This survey (re)introduces reinforcement learning to researchers studying sequential decision problems. I review the theoretical connections between dynamic programming and reinforcement learning, demonstrating how value iteration, Q-learning, and policy gradient methods are common solution methods to the same class of optimization problems. I then examine applications across several domains, including structural economic models with high-dimensional state spaces; strategic games in which multi-agent algorithms compute equilibria under imperfect information; bandit problems in which economic structure yields tighter regret bounds; preference learning; and field deployments in recommendations, marketplaces, operations, physical systems, and post-training. The exposition combines formal theory, practical applications and computational illustration.

Both dynamic programming and reinforcement learning solve the Bellman equation; they differ in the information requirements and the way in which the solution is refined. First, dynamic programming requires knowledge of the transition in the environment and the reward function which allows the reduction of the average Bellman error, reinforcement learning estimates value functions only from sampled transitions (observed sets of state, action, reward, next-state) which only allows reduction of the sampled Bellman error \textit{at that state-action pair}. This allows us to improve policies in domains where it is easier to build a simulator than specify the model of the environment and rewards e.g. board games, physics simulators for robots. Second, dynamic programming makes a ``breadth-first'' (across all states and actions) update of the solution at each sweep, while reinforcement learning makes a ``incremental'' (only for the current state and action) update; this greatly reduces the computational burden and enhances scalability.

Reduction of average Bellman errors gives dynamic programming a geometric rate of convergence to the optimal solution, while the incremental updates and reduction of sampled Bellman errors, when combined with ``sufficient exploration'' of the state-action space, gives reinforcement learning only sublinear convergence guarantees. The trade is favourable where exact dynamic programming cannot be computed at all; where it can, the simulations in this survey generally find it matching or beating the sampled alternative. These approximation methods sacrifice theoretical guarantees. RL algorithms lack the convergence assurances of exact dynamic programming. They exhibit sensitivity to hyperparameters and initialization. They can converge to suboptimal policies without diagnostic indication. This survey presents reinforcement learning as a computationally flexible framework while acknowledging its methodological limitations, which Section~\ref{section:deeprl_practice} documents in detail.

Theory in reinforcement learning trails empirical success, often by years; convergence guarantees, sample complexity bounds, and approximation error characterizations typically arrive after practitioners have demonstrated that an algorithm works. The theoretical insights that eventually follow tend to be deep and structural, and the empirical frontier is itself a productive research frontier. Experiments are brittle, conducted on benchmark environments that are stylized approximations of deployment settings. These benchmarks nonetheless serve a critical coordination function, aligning research effort, enabling reproducible comparison, and exposing failure modes that motivate new theory. Details matter disproportionately in reinforcement learning; small implementation choices can determine whether an algorithm converges or diverges, and the practice of releasing code and documenting hyperparameters, seeds, and preprocessing has proven essential to progress.

Existing surveys of reinforcement learning for economists divide along scope. \citet{CharpentierElieRemlinger2021} present reinforcement learning techniques alongside applications in economics, game theory, operations research and finance, and treat inverse reinforcement learning and its dynamic discrete choice lineage. \citet{MosaviEtAl2020} review deep learning and deep reinforcement learning methods, with the reinforcement learning applications concentrated in stock trading. \citet{HamblyXuYang2023} work through the finance applications, namely optimal execution, portfolio optimization, option pricing and hedging, market making, robo-advising and order routing. \citet{AtashbarShi2022} cover deep reinforcement learning for macroeconomics, organized around solution methods and bounded rationality. \citet{iskhakov2021} contrast machine learning with structural econometrics, asking where machine learning can advance the goals of structural work, and set inverse reinforcement learning against structural estimation.

This survey differs in three respects. It organizes the material around the Bellman equation and the convergence of its solution methods rather than around application domains. It develops material that the earlier surveys do not, including offline reinforcement learning, learning from human feedback, and distributionally robust and constrained methods. Each chapter carries a working simulation benchmarked against a dynamic programming or analytical solution.

Several topics are treated elsewhere. Algorithmic collusion, in which independent pricing algorithms learn to sustain supra-competitive prices \citep{Calvano2020}, is treated in a companion thesis chapter \citep{Rawat2026collusion} and omitted here. Portfolio optimization, optimal execution, and asset pricing via reinforcement learning form a large body of work surveyed comprehensively elsewhere \citep{HamblyXuYang2023}. The inverse problem of inferring preferences from observed behavior using inverse reinforcement learning and structural estimation is treated in a companion survey \citep{RustRawat2026}.

The survey addresses the forward problem, that is, computing optimal policies given a known or simulated environment. Section~\ref{section:history} traces the parallel historical development of dynamic programming and reinforcement learning. Section~\ref{section:rl_algorithms} surveys reinforcement learning algorithms, and Section~\ref{sec:planning_learning} develops the unified theory connecting planning and learning, with Section~\ref{section:deeprl_practice} examining the empirics of deep RL. Section~\ref{section:rl_econ_models} applies reinforcement learning to structural estimation of econometric models. Section~\ref{section:rl_macro} treats macroeconomic models, and Section~\ref{section:rl_games} addresses strategic games. Section~\ref{section:bandits} surveys bandit problems with domain structure. Section~\ref{section:offline_rl} covers offline reinforcement learning and Section~\ref{section:rlhf} covers reinforcement learning from human feedback. Section~\ref{section:causal_rl} treats causal inference for RL, with Section~\ref{section:rl_for_ci} treating RL methods for causal inference. Section~\ref{section:dist_robust_constrained} surveys quantile, robust, and constrained reinforcement learning, Section~\ref{section:world_models} treats world models and model-based RL, and Section~\ref{section:field_deployments} synthesizes the narrower evidence on real-world RL deployments. Section~\ref{section:conclusion} concludes.
